% Generated draft: requires mathematical and source review.
\documentclass[11pt]{article}
\usepackage[margin=25mm]{geometry}
\usepackage{fontspec}
\setmainfont{Latin Modern Roman}
\usepackage{amsmath,amssymb,mathtools,unicode-math}
\setmathfont{Latin Modern Math}
\usepackage{enumitem,needspace,xurl,hyperref}
\setlength{\parindent}{0pt}
\setlength{\parskip}{5pt}
\setlength{\emergencystretch}{3em}
\newcommand{\N}{\mathbb{N}}
\newcommand{\Z}{\mathbb{Z}}
\newcommand{\Q}{\mathbb{Q}}
\newcommand{\R}{\mathbb{R}}
\newcommand{\C}{\mathbb{C}}
\newcommand{\F}{\mathbb{F}}
\newcommand{\A}{\mathbb{A}}
\newcommand{\PP}{\mathbb{P}}
\newcommand{\E}{\mathbb{E}}
\newcommand{\eps}{\varepsilon}
\newcommand{\dd}{\,\mathrm{d}}
\begin{document}
\section*{UnsolvedMath ID: 30005924}
\textbf{Source problem code:} catalogue entry
\subsection{Definitions and mathematical statement}
A noncommutative polynomial \(Q\) in indeterminates \(t_1,\ldots,t_n\) is a finite linear combination of words in those indeterminates. The supplied record does not specify its coefficient field, an evaluation tuple, or the ambient operator algebra.

For a Hilbert space \(H\), an operator \(T:H\to H\) has finite rank if its range is finite dimensional. A bounded operator \(T:H\to H\) is Hilbert--Schmidt if \(\sum_{e\in\mathcal B}\|Te\|^2<\infty\) for an orthonormal basis \(\mathcal B\) of \(H\). These definitions explain the two operator classes mentioned in the source; the source excerpt does not identify the relevant Hilbert space or operators.

The supplied question asks what regularity conclusions can be obtained under various regularity assumptions on free conjugate variables. It also asks when certain \(1\)-cycles relevant to \(L^2\)-homology can have components \(T_j\) that are finite-rank operators, as compared with Hilbert--Schmidt operators. The excerpt does not specify the underlying noncommutative probability space, the variables and conjugate-variable relations, the regularity hypotheses or conclusions, the chain complex, or the cycle equation. Consequently, it does not determine a precise quantified mathematical problem.

The excerpt additionally prints the formal relation
\[
Q\otimes 1-1\otimes Q
=\sum_{i=1}^{n}\bigl(t_i\partial_iQ-(\partial_iQ)t_i\bigr),
\]
where \(\partial_i\) are called noncommutative difference quotients. The tensor-space conventions and the actions denoted by multiplication with \(t_i\) are not supplied. This formula is retained as a source fragment, without choosing multiplication conventions or treating it as the unresolved claim.

The supplied statement-verification field says that the report uses the broad question to organize a survey of established regularity results and does not isolate an unresolved claim. Its statement status is \texttt{unrecoverable}. Thus no definite conjecture can be recovered from this record without additional source context.
\subsection{Short English statement}
The record asks broadly how assumptions on free conjugate variables affect regularity and mentions finite-rank versus Hilbert--Schmidt components of \(L^2\)-homology cycles. Essential hypotheses and the intended conclusion are missing. The supplied verification describes a survey question rather than an isolated unresolved claim, so this draft preserves the fragment and its limitations.
\subsection{Sources}
\begin{itemize}
\item Local UnsolvedMath registry: supplied complete record, permanent ID \(30005924\), title ``Free-Conjugate-Variable Regularity Overview,'' exact source problem code \texttt{catalogue entry}. The record supplies proposed year \(2024\) and no named proposer.
\item Supplied source citation: \emph{Non-commutative Function Theory and Free Probability} (2024), Oberwolfach Reports, pp.\ 1227--1284, DOI \texttt{10.4171/owr/2024/22}. \url{https://doi.org/10.4171/owr/2024/22}. The supplied background describes Oberwolfach Reports as documenting mathematical workshops and their open-problem discussions.
\item Supplied literature source ``OWR 2024/22'': \url{https://doi.org/10.4171/owr/2024/22}.
\item Supplied literature source ``Regularity results'': \url{https://par.nsf.gov/servlets/purl/10432807}. No author, title, or specific theorem is supplied for this reference.
\item Supplied literature assessment dated exactly \texttt{2026-08-22}: the cited free-probability results provide regularity and homological information in important cases, while the broad regularity question is described as remaining unresolved. This is supplied dated curation derived from the OWR record, not an independently verified literature assessment; the background expressly requests independent verification of time-sensitive status and cited claims.
\item The registry labels the entry \texttt{open}, but also labels its statement \texttt{unrecoverable}. Its supplied statement verification says the official report organizes a survey of established regularity results around this broad question and does not isolate an unresolved claim. These are supplied classifications and qualifications, not independently verified mathematical facts.
\end{itemize}
\end{document}
